\documentclass[11pt,a4paper]{article}
\usepackage[T1]{fontenc}
\usepackage{lmodern,amsmath,amssymb,amsthm,booktabs,longtable,array,geometry,microtype}
\usepackage[hidelinks]{hyperref}
\geometry{margin=26mm}
\newtheorem{theorem}{Theorem}
\newtheorem{lemma}[theorem]{Lemma}
\newtheorem{corollary}[theorem]{Corollary}
\newtheorem{proposition}[theorem]{Proposition}
\newcommand{\Ber}{\operatorname{Ber}}
\newcommand{\Prob}{\mathbb P}
\newcommand{\R}{\mathrm R}
\newcommand{\U}{\mathrm U}
\title{An exact all-parameter verification of the bump--unwind\\
path order in Bernoulli last-passage percolation for $n=4$}
\author{Unsigned research draft}
\date{29 October 2026\\\small Research manuscript; not submitted to a journal; not peer reviewed}
\begin{document}
\maketitle
\begin{abstract}
We verify a finite specialization of the path-probability ordering proposed
by McKeown and Nakajima for vertex last-passage percolation. On the
$5\times5$ vertex grid with independent Bernoulli($p$) weights, every
genuine nonidentity bump or unwinding move strictly increases the
probability of membership in the full set of maximizing paths for
$0<p<1$. Equality holds at $p=0,1$. The assertion covers all 120
single-step comparisons and their 188 distinct ordered pairs in the
transitive closure. Two independently written enumerators agree on every
integer coefficient obtained from all $2^{23}$ free-vertex configurations.
For 118 edges, positivity follows directly from their degree-23
nonnegative coefficient arrays. The two remaining edges have a negative
coefficient, but a single exact degree elevation yields nonnegative
coefficients for every edge. The exceptional edges also disprove a
stronger comparison conditional on exactly three occupied free vertices.
We document boundary conventions, complete certificates, small-grid hand
checks, and reproducibility. The result is a modest computer-assisted
finite verification; priority and independent research significance are
not established.
\end{abstract}

\section{Scope and relation to prior work}
The question of which directed paths have greatest probability of being
maximal has a geometric history; see Alberts and Cator \cite[Section~8]{AC}
for modal-path and path-order questions. McKeown and Nakajima \cite{MN}
give two concrete path transformations and propose an increase of
geodesic-membership probabilities. Their Theorem~3, equation~(6), proves
the Bernoulli comparison for sufficiently small positive $p$, with
strictness for distinct related paths. The choice in their Section~5 is
\[
p_n=\min\left\{\frac1{2N},\frac1{8\binom N3}\right\},
\qquad N=(n+1)^2,
\]
which gives $p_4=1/18400$. Our computation covers the whole parameter
interval for the fixed size $n=4$, and checks $n\leq3$ as controls.
It neither proves their general conjecture nor concerns the separate
moderate-deviation issue resolved in their version~2.

We use version~2, dated 20 September 2026; the version record was checked
on 29 October 2026. A bounded search of 15 query strings, author materials,
and relevant backward references did not locate the same finite
certificate. Forward citation coverage is incomplete: Google Scholar,
ADS, and Semantic Scholar could not be fully inspected, and one author's
homepage was inaccessible. Thus no priority assertion follows. The
full query log, URLs, access failures, and distinctions are in
\path{literature/audit.txt} and \path{literature/sources.json}.
This unsigned manuscript is prepared for public archival on GitHub.
It has not been submitted to a journal or sent to the cited authors.
A publication-stage refresh with ten additional queries is recorded in
\path{literature/novelty_check_20261029.txt}; priority remains unestablished.

\section{Model and exact meaning of the path order}
Let $V_n=\{0,\ldots,n\}^2$ and let $\Gamma_n$ consist of the $\binom{2n}n$
words containing $n$ letters $\R$ and $n$ letters $\U$. A word defines the
corresponding up/right path from $(0,0)$ to $(n,n)$.
For independent weights $\omega_v\sim\Ber(p)$ put
\[
L(\gamma)=\sum_{v\in\gamma}\omega_v,\qquad
L_n=\max_{\eta\in\Gamma_n}L(\eta),\qquad
F_\gamma(p)=\Prob_p\{L(\gamma)=L_n\}.
\]
Every tied maximizing path is counted. In particular, these probabilities
do not in general sum to one; we do not select one path from a tie.

\subsection{Run notation and bump moves}
Write a word as its maximal constant-letter runs with signed lengths
$a=(a_1,\ldots,a_r)$, positive for $\U$ and negative for $\R$.
The nonzero signs alternate, the positive lengths sum to $n$, and the
absolute negative lengths sum to $n$. Let $s_i=\sum_{j=1}^i a_j$ and
$a_0=s_0=0$. Thus $s_i$ is the vertical coordinate minus the horizontal
coordinate at the end of run $i$.

A bump has $2\leq k\leq r-2$ and satisfies either
$s_k\geq0$, $a_k<0$, or $s_k\leq0$, $a_k>0$. Its replacement is
\[
(a_1,\ldots,a_{k-2},\
 a_{k-1}+a_{k+1},\ a_k+a_{k+2},\
 a_{k+3},\ldots,a_r).
\]
Equivalently, reverse the order of the letters between corner times
$t_{k-1}$ and $t_{k+1}$, where
$t_i=\sum_{j=1}^i|a_j|$. The blocks on either side merge with the resulting
blocks, reducing the number of runs by two.
The weak inequalities permit diagonal contact. The putative index
$k=1$ cannot satisfy the sign condition because $s_1=a_1$.

\subsection{Unwinding moves and degenerate indices}
A genuine unwind has $0\leq l$, $l+2\leq k\leq r-1$, and a sign
$\varepsilon\in\{-1,1\}$ such that
\[
 \varepsilon s_k>0,\qquad
 \varepsilon s_i<0\quad(l<i<k),\qquad
 l=0\ \text{or}\ \varepsilon s_l>0.
\]
Reverse the order of all letters between times $t_l$ and $t_k$, keeping
each letter's direction. For $l\geq1$ the resulting signed runs are
\[
(a_1,\ldots,a_{l-1},\
a_l+a_k,\ a_{k-1},\ldots,a_{l+2},\
a_{k+1}+a_{l+1},\ a_{k+2},\ldots,a_r).
\]
For $l=0$, the prefix is absent and the first displayed merged term
is simply $a_k$. This case reverses an initial segment, and decreases
the number of runs by one; an interior unwind decreases it by two.

These conventions explicitly reconcile the run display with the
geometric operations in \cite[Section~2, Figures~3--4]{MN}. The strict
interior signs express an actual excursion to the other side of the
diagonal. If one allows $k=l+1$ in the prose inequality, its interior
condition is vacuous; the geometric segment then consists of a single
constant-letter run and its reversal is an identity. The merged-run
display must not be applied with overlapping indices in that case.
It is excluded from our nonidentity move list. This is a stated
interpretation of the genuine geometric moves, not a claim about a
literal overlapping-index substitution. The source's assertion that
$n=1$ has no nontrivial moves and its Lemma~9 strict two-dot gain agree
with this interpretation.

The initial boundary $l=0$ is expressly retained. There is no extra
terminal boundary: $s_r=0$ makes $k=r$ fail the strict condition, and
$a_{k+1}$ consequently exists. Diagonal contacts at interior corners
fail the strict unwind condition. No additional moves are introduced
by reversing a whole path. Diagonal reflection, on the other hand,
is exactly the sign symmetry already included in the definitions.
An independent implementation operates on corner coordinates and word
segments rather than on the merged-run formulas.

Write $\gamma\prec\pi$ when $\pi$ is reachable from $\gamma$ by one or
more such moves. The run count strictly decreases, so this is a strict
partial order. The identity belongs only to its reflexive closure.

\section{Finite theorem}
\begin{theorem}[Computer-assisted finite comparison]\label{thm:main}
For $1\leq n\leq4$, and every genuine nonidentity single-step move
$\gamma\longrightarrow\pi$ defined above,
\[
 F_\pi(p)>F_\gamma(p)\qquad(0<p<1).
\]
The same strict inequality holds if $\gamma\prec\pi$.
For $p=0$ or $p=1$, every path is maximizing, so $F_\gamma(p)=1$.
For $n=4$, each single-step difference satisfies the explicit bound
\[
 F_\pi(p)-F_\gamma(p)\ \geq\ 2p^2(1-p)^{22}
 \qquad(0\leq p\leq1).
\]
\end{theorem}
\begin{corollary}
For $n=4$ and $0<p<1$, the only modal paths are
$\R^4\U^4$ and $\U^4\R^4$. Their probabilities are equal, and exceed
that of each of the other 68 paths.
\end{corollary}
The corner statement follows because every noncorner reaches a corner
and the two corners have equal probability under diagonal reflection.
This does not compare every incomparable pair in a total order.

\section{Exact reduction and sign certificates}
\begin{lemma}[Removal of common endpoints]\label{lem:endpoints}
Setting the weights at $(0,0)$ and $(n,n)$ to zero preserves the maximizing
set for every configuration.
\end{lemma}
\begin{proof}
Each path contains both vertices, so the same sum of two weights is
subtracted from every score. All equalities and inequalities are preserved.
\end{proof}

Put $W_n=V_n\setminus\{(0,0),(n,n)\}$ and $m=|W_n|=(n+1)^2-2$.
For $S\subseteq W_n$ define
\[
 M(S)=\max_{\eta\in\Gamma_n}|S\cap\eta|,\quad
 C_\gamma[k]=
 \#\{S\subseteq W_n:|S|=k,\ |S\cap\gamma|=M(S)\}.
\]
Every set $S$ is equally likely within its cardinality layer, and its
unconditional probability is $p^k(1-p)^{m-k}$. Hence, with $q=1-p$,
\begin{equation}\label{eq:raw}
 F_\gamma(p)=\sum_{k=0}^m C_\gamma[k]p^kq^{m-k}.
\end{equation}
The $C_\gamma[k]$ are raw configuration counts. There is no extra
$\binom mk$ multiplier. If one uses the normalized Bernstein basis
$\binom mk p^kq^{m-k}$, its coefficient is $C_\gamma[k]/\binom mk$.

For an edge write $d_k=C_\pi[k]-C_\gamma[k]$. When all $d_k\geq0$ and
some $d_k>0$, equation~\eqref{eq:raw} gives strict positivity for
$0<p<1$. A negative $d_k$ does not by itself refute the comparison.

\begin{lemma}[Exact degree elevation]\label{lem:elevation}
Set $d_{-1}=d_{m+1}=0$ and $e_k=d_k+d_{k-1}$ for $0\leq k\leq m+1$.
Then
\[
 \sum_{k=0}^m d_kp^kq^{m-k}
 =\sum_{k=0}^{m+1}e_kp^kq^{m+1-k}.
\]
If all $e_k\geq0$ and $e_2\geq2$, the difference is at least
$2p^2q^{m-1}$.
\end{lemma}
\begin{proof}
Multiply the first sum by $p+q=1$ and collect terms of the same
degree. All basis terms are nonnegative on $[0,1]$ and positive in
its interior, proving the claimed bound.
\end{proof}

For all 120 edges at $n=4$, the saved arrays satisfy $e_k\geq0$ and
$e_2\geq2$. The checks are exact integer comparisons. The same condition
holds for every edge at $n=2,3$; $n=1$ has no edges.
No root solver, floating-point sign test, or numerical grid in $p$
is used.

\subsection{The two exceptional degree-23 arrays}
The only edges failing raw coefficientwise positivity at $n=4$ are
\[
 \R\U\U\R\U\U\R\R\longrightarrow\U\U\R\R\U\U\R\R
 \quad\text{and its diagonal reflection.}
\]
The first word has runs $(-1,2,-1,2,-2)$ and heights
$s_1=-1,s_2=1$. Thus $l=0,k=2$ is legal: reversing the prefix
$\R\U\U$ gives $\U\U\R$. Its destination has runs $(2,-2,2,-2)$.
Both edges have the identical difference array in Table~\ref{tab:exception}.

\begin{table}[ht]
\centering
\begin{tabular}{r r r @{\qquad} r r r}
\toprule
$k$&$d_k$&$e_k$&$k$&$d_k$&$e_k$\\
\midrule
0&0&0&13&2596&6225\\
1&0&0&14&1347&3943\\
2&2&2&15&536&1883\\
3&-2&0&16&204&740\\
4&40&38&17&94&298\\
5&67&107&18&44&138\\
6&328&395&19&12&56\\
7&864&1192&20&1&13\\
8&1530&2394&21&0&1\\
9&2832&4362&22&0&0\\
10&3783&6615&23&0&0\\
11&3948&7731&24&0$^\dagger$&0\\
12&3629&7577&&&\\
\bottomrule
\end{tabular}
\caption{Exceptional edge coefficients. The left representation has degree
23 and the elevated representation degree 24.
$^\dagger d_{24}=0$ is only a padding convention.}\label{tab:exception}
\end{table}
Here $d_3=-2$ is offset exactly by $d_2=2$ upon elevation; the next
coefficient becomes $e_4=38$. As a second elementary sign check, the
first three nonzero-degree terms can be grouped as
\[
p^2q^{19}(2q^2-2pq+40p^2)
=\tfrac12p^2q^{19}\bigl((2q-p)^2+79p^2\bigr),
\]
and every remaining nonzero term is positive in the interior.

\begin{proposition}[Failure of a stronger fixed-cardinality comparison]
Conditioning on exactly three occupied vertices in $W_4$, the preceding
legal unwind changes geodesic-membership probability by $-2/1771$.
Thus monotonicity in each fixed-cardinality ensemble is false.
\end{proposition}
\begin{proof}
The conditional difference equals $d_3/\binom{23}3=-2/1771$.
This conditional model is different from the independent Bernoulli model;
the proposition is not a counterexample to Theorem~\ref{thm:main}.
\end{proof}

\section{Complete enumeration and verification}
\subsection{Two algorithms}
Free vertices are indexed in increasing $y$, then increasing $x$,
omitting the two endpoints. Paths are indexed by their complete words
in lexicographic order $\R<\U$.

Implementation A, \path{src/count_masks.cjs}, constructs a bit mask
$B_\gamma$ for each path's free vertices. For each integer
$0\leq b<2^m$, it computes every score
$\operatorname{popcount}(b\mathbin{\&}B_\gamma)$, takes their maximum,
and increments the count of \emph{every} path attaining it in the
$\operatorname{popcount}(b)$ column. Its population count uses fixed
integer bit operations.

Implementation B, \path{src/replay_dp.js}, independently generates paths
as lists of grid indices and unpacks each configuration into individual
cell weights. It obtains the global maximum by the lattice recurrence
\[
 D(x,y)=\omega_{x,y}+\max\{D(x-1,y),D(x,y-1)\},
\]
using only the available predecessor on each boundary and
$D(0,0)=0$. Every path's weights are then added directly and compared
with $D(n,n)$. This implementation was written before inspecting A's
scoring code; it neither imports A's masks nor obtains the maximum from
its path scores. The shared items are the mathematical model and
indexing specification.

\begin{lemma}[Enumeration correctness]
Each algorithm computes the integers in equation~\eqref{eq:raw}.
\end{lemma}
\begin{proof}
The mask loop is a bijection from the integers $0,\ldots,2^m-1$ to
subsets of the free vertices. Recursive word generation lists every
legal path once. In A, intersection population counts equal occupied
path-vertex counts, and maximization is over the complete path set.
In B, induction on $x+y$ proves that $D(x,y)$ is the greatest path
weight to $(x,y)$; a path's separately summed score equals $D(n,n)$
if and only if it is maximal. Both algorithms increment precisely
the required integer once for every successful pair $(S,\gamma)$.
\end{proof}

At $n=4$ the maximum score is seven, and every count is at most
$2^{23}$. The unsigned 8-bit scores and 32-bit counters are sufficient.
All relevant JavaScript integers are exact. In A the largest product
inside the population-count reduction is below $2^{53}$, so the
floating-number representation does not affect its integer result.
No probabilities are accumulated in floating point.

\subsection{Coverage and invariants}
\begin{table}[ht]
\centering
\begin{tabular}{r r r r r r r}
\toprule
$n$&$m$&Configurations&Paths&Bumps&Unwinds&Closure pairs\\
\midrule
1&2&4&2&0&0&0\\
2&7&128&6&2&2&4\\
3&14&16,384&20&14&10&30\\
4&23&8,388,608&70&76&44&188\\
\bottomrule
\end{tabular}
\caption{Complete coverage. Closure pairs exclude identities. At $n=4$,
30 of the 44 unwinds use $l=0$. Each listed operation has a distinct
ordered source--target pair.}
\end{table}
Both count implementations agree on all $6+48+300+1680=2034$ entries
for $n=1,2,3,4$. They separately check the cardinality layers against
$\binom mk$, the endpoint columns $C_\gamma[0]=C_\gamma[m]=1$,
and $0\leq C_\gamma[k]\leq\binom mk$. Counts agree under both diagonal
reflection and central rotation with reversal of traversal (word reversal).
These symmetries are checks on probabilities, not a license to add moves.

The signed-run and geometric word implementations agree on every
operation witness $(\gamma,\pi,\mathrm{kind},l,k)$, not just on edge
totals. Every resulting word is valid and its run count decreases by
the stated amount. Their complete transitive closures agree, and the
only sinks are the two corner paths.

A further check directly counts incomparable vertex pairs. Let
\[
 I(\gamma)=
 \#\{(v,w):v\in\gamma,\ w\in W_n\setminus\gamma,\
                   (v_x-w_x)(v_y-w_y)<0\}.
\]
Then
\[
 C_\gamma[1]=2n-1,\qquad
 C_\gamma[2]=\binom{2n-1}{2}+I(\gamma).
\]
Indeed, either both occupied sites lie on $\gamma$, or exactly one
does and the two sites cannot lie together on any directed path.
If neither lies on $\gamma$, that path cannot be maximizing.
The incomparable-pair calculation matches every coefficient and yields
$I(\pi)-I(\gamma)>0$ on every edge. This independently checks the
two-dot strictness underlying \cite[Lemma~9]{MN} without relying on
its area calculation.

\subsection{Small cases by hand}
For $n=1$ each path has one nonendpoint weight, and success is the event
$X\geq Y$. Its probability is $1-p+p^2$, with raw count row $(1,1,1)$.
There is no genuine move.

For $n=2$, label the seven free weights
\[
a=\omega_{1,0},\ b=\omega_{2,0},\ c=\omega_{0,1},\
d=\omega_{1,1},\ e=\omega_{2,1},\ f=\omega_{0,2},\ g=\omega_{1,2}.
\]
The path scores in order
$\R\R\U\U,\R\U\R\U,\R\U\U\R,\U\R\R\U,\U\R\U\R,\U\U\R\R$
are
\[
 a+b+e,\quad a+d+e,\quad a+d+g,\quad
 c+d+e,\quad c+d+g,\quad c+f+g.
\]
Let $G_\gamma(t)=\sum_k C_\gamma[k]t^k$,
$Q=1+t+t^2$, and $E=1+t^2$. The three symmetry classes have
\begin{align*}
G_{\R\U\R\U}&=(1+t)Q^3-tE^2,\\
G_{\R\U\U\R}&=t(1+t)^2Q^2+(1+t+2t^2)^2,\\
G_{\R\R\U\U}&=(1+t)(1+t^2)Q^2-tE^2+tH,\\
H&=1+5t+7t^2+9t^3+3t^4+t^5.
\end{align*}
For the first class, the score comparisons require $d\geq b$,
$e\geq g$, and $a\geq c$. The sole remaining failure has
$d=b=0$, $a=c$, $e=g$, $f=1$, giving the subtraction $tE^2$.
For the second, $a\geq c$, $g\geq e$, and
$d+(g-e)\geq b$, $d+(a-c)\geq f$ suffice; splitting on $d$ gives
the two summands. For the corner, split the constraint $b\geq d$
into $(b,d)=(0,0),(1,1),(1,0)$. The first two cases give the first
two terms combined above. In the third case the remaining comparison
is $c+f+g\leq a+e+1$, yielding
\[
H=\sum_{A=0}^2\binom2A t^A
       \sum_{B=0}^{\min(3,A+1)}\binom3B t^B.
\]
Expansion gives
\[
\begin{array}{c|rrrrrrrr}
\gamma\backslash k&0&1&2&3&4&5&6&7\\\hline
\R\R\U\U&1&3&11&13&17&8&4&1\\
\R\U\R\U&1&3&9&11&13&8&4&1\\
\R\U\U\R&1&3&9&12&14&8&4&1
\end{array}
\]
and symmetry gives the remaining rows. The differences from the corner
are $2t^2+2t^3+4t^4$ and $2t^2+t^3+3t^4$. Thus all 48 small-case
entries have a separate algebraic check.

\subsection{Completion of the finite proof}
The complete counts, exact move coverage, and sign certificates establish
every single-step assertion of Theorem~\ref{thm:main} by
Lemmas~\ref{lem:endpoints} and \ref{lem:elevation}. Summing the positive
single-step differences along a nonempty chain establishes the strict
partial-order assertion. Direct evaluation of the deterministic all-zero
and all-one environments gives the endpoint equalities.

The finite data, rather than an unreported search, are part of this
computer-assisted proof. The independent replay is complete, not
statistical. The proof still relies on correct execution of the supplied
programs and ordinary arithmetic implementations; it is not a
proof-assistant formalization or a human peer-reviewed theorem.

\section{Reproducibility, audit, and limitations}
The project includes:
\begin{itemize}
\item \path{results/counts_n4.csv} and the corresponding JSON: all
70 count rows, with 24 entries each;
\item \path{results/moves_n4.json}: all 120 legal operations and witnesses;
\item \path{results/certificates_n4.json}: every difference array, its
degree-24 array, and all 188 strict closure pairs;
\item independent count and move outputs, and reports in \path{review/};
\item \path{src/certify.py}, \path{reproduce.py}, and
\path{README.md}: exact validation and complete regeneration;
\item \path{literature/}: source/version records and bounded search log;
\item \path{SHA256SUMS.txt}: integrity hashes for retained artifacts.
\end{itemize}
No nonstandard Python package is required. The count implementations
use Node.js; the checker uses Python's standard library.
The initial $n=4$ runs used Node v24.19.0 on Windows x64:
A took 1.377 wall seconds and 1.390 CPU seconds; B took 4.563 wall seconds
and 4.500 CPU seconds, with approximately 40.6 MiB peak resident memory
for B. Short benchmarks preceded full runs. The initial full runs used
5.890 CPU seconds combined, well inside the 7200-second limit.
The supplied portable reproduction script runs workers sequentially and rejects
$n>4$. Timings are observations on this machine, not portable guarantees.
Rerun timing receipts have different hashes; deterministic count arrays
and certificate arrays are the mathematical comparison targets.

Independent AI agents audited move definitions, rebuilt the full count
table, checked the two exceptional sign certificates, and reviewed the
manuscript. These activities are distinct from human expert review.
There is no conclusion for $n\geq5$, no assertion about other weight
laws, and no claim that the finite computation is new globally or
has standalone publication significance. The exceptional negative
fixed-cardinality coefficient explains why positivity of every original
coefficient is too strong a proposed general proof strategy; one degree
elevation happens to suffice here.

\begin{thebibliography}{9}
\bibitem{MN}
S. McKeown and S. Nakajima,
\emph{On the most likely geodesic in last passage percolation},
arXiv:2608.11378v2, 20 September 2026.
\url{https://arxiv.org/html/2608.11378v2}.
Version record: \url{https://arxiv.org/abs/2608.11378}.
Accessed 29 October 2026.

\bibitem{AC}
T. Alberts and E. Cator,
\emph{On the Passage Time Geometry of the Last Passage Percolation Problem},
ALEA, Latin American Journal of Probability and Mathematical Statistics
\textbf{18} (2021), 211--247.
doi:10.30757/ALEA.v18-10.
\url{https://alea.impa.br/articles/v18/18-10.pdf}.
Relevant comparison: Section~8, pp.~242--244.
Accessed 29 October 2026.
\end{thebibliography}
\end{document}
